% ==========================================================
% 040_baseline_error_anatomy.tex
% Forecast Admissibility Surface (FAS)
% ==========================================================

\section{Baseline Error Anatomy}
\label{sec:fas_baseline_error_anatomy}

The Forecast Admissibility Surface is derived from observable properties of
baseline forecast error rather than from the performance of advanced or
optimized models. This design choice is intentional. Baseline models provide a
low-capacity, transparent reference that exposes structural characteristics of
a demand slice without masking them through complex fitting or feature
interaction. FAS relies on these baseline-derived signals to assess whether a
slice is amenable to learning and control at all.

Let $\yit$ denote realized demand and $\yhatit$ denote the corresponding baseline
forecast for entity $\entity$ at time index $\timeindex$. FAS operates on the
distribution of absolute error
$\absit = \left|\yhatit - \yit\right|$ aggregated over each slice in the chosen
domain. The resulting \emph{error anatomy} summarizes both the magnitude and the
structure of baseline forecast deviations.

\subsection{Support Size}

The most fundamental prerequisite for admissibility is sufficient observational
support. For each slice, FAS records the number of observations $n$ available for
baseline error analysis. Slices with extremely limited support cannot be
reliably characterized: apparent stability or instability may be an artifact of
sampling rather than structure.

Rather than attempting to infer hostility from sparse data, FAS treats low
support conservatively. Slices failing to meet a minimum support threshold are
classified as \Conditional, reflecting epistemic uncertainty rather than
demonstrated inadmissibility. This prevents the exclusion of slices solely due to
data scarcity while avoiding unwarranted confidence.

\subsection{Spike Frequency}

Operational demand often exhibits rare but severe deviations from baseline
behavior, driven by exogenous shocks, executional anomalies, or bursty demand
patterns. To capture this behavior, FAS computes the \emph{spike rate}: the
fraction of observations in a slice for which absolute error exceeds a
predefined spike threshold.

High spike rates indicate that baseline forecasts routinely fail by large
margins, even if average error appears moderate. Such slices are particularly
hostile to stable learning, as infrequent extremes dominate both training loss
and evaluation metrics. Spike frequency therefore serves as a primary signal of
structural instability.

\subsection{Tail Magnitude}

In addition to spike frequency, FAS characterizes the severity of extreme errors
via upper-tail quantiles of the absolute error distribution. Specifically, high
quantiles (e.g., the $90^{\text{th}}$ or $95^{\text{th}}$ percentile) summarize
the magnitude of errors encountered in the worst-performing regimes of a slice.

Elevated tail magnitude signals that even when extreme events are rare, their
impact is substantial. In such cases, forecasts may appear acceptable under
average metrics while remaining operationally unsafe. Tail-based diagnostics
therefore complement spike frequency by distinguishing between frequent moderate
instability and infrequent but catastrophic failure.

\subsection{Baseline-Centric Interpretation}

It is important to emphasize that baseline error anatomy is not intended as a
proxy for final forecast quality. A slice may exhibit poor baseline performance
yet still be improved through modeling, while another may resist improvement
entirely. FAS does not attempt to predict which outcome will occur. Instead, it
uses baseline behavior as a \emph{structural probe}: if instability is already
severe at low model capacity, the slice warrants heightened scrutiny or
restriction before additional complexity is introduced.

By grounding admissibility in baseline error anatomy, FAS ensures that decisions
about where forecasting is allowed are driven by observable, interpretable
signals rather than by optimistic assumptions about model capability.
